\section{All charges of the original crossing traces}\label{ms:filling}
The powers $n_v$ are fixed signed integers. In the primitive filling
convention they are $n_v=\epsilon_v-2p_v$; odd powers are included.
This section proves Theorem~\ref{ms:filled-main} analytically. It
does not identify a geometric action.

\subsection{The full finite functional and its currents}
For an even polynomial $f(\tau)=f^\sharp(-2\cos\tau)$ define
$\mathcal D_{n,\sigma,N}f=\phi_\sigma(R_\sigma^n f^\sharp)$.
The full crossing data in \cite[Proposition 6.1 and equation (6.4)]{Meridian}
give, with $\eta=(-1)^d$,
\begin{align}
 \mathcal D_{n,\sigma,N}f
 &=\frac{\eta N C_{N,n}(\sigma)}{\sin2\pi\sigma}
 \sum_{r=0}^{N-1}\sin\tau_r\,f(\tau_r)
                  e^{\ii nN\tau_r^2/(4\pi)},\label{ms:trace}\\
 \tau_r&=\frac{2\pi}{N}(d/2+\sigma-r),\qquad
 C_{N,n}(\sigma)=q^{n(d^2/4-\sigma^2/2)}.
\end{align}
This is first written off the collisions. Its removable limit is
the full Hermite functional, not just the node values. In
particular inverse powers use the same formula.

Introduce independent formal grid parameters $\sigma_v$ temporarily;
the actual assignments are $0,s,-s$. If $\zeta_N$ is the original
closing and framing scalar and $\rho_{v,N}$ the reference scalar
for its orientation type, set $\widehat\zeta_N=\zeta_N\prod_v\rho_{v,N}$.
The exact source reduction already proved has the form
\begin{equation}\label{ms:source-to-functional}
 F_N(q^s)=\widehat\zeta_N(s)H_N(s)
       \left(\prod_v\mathcal D_{n_v,\sigma_v,N}^{(v)}\right)
       B_N(\boldsymbol\tau,s)
       +O(N^A e^{N(F_0-\eta_0)}).
\end{equation}
Each factor acts in its own variable; this is tensor evaluation,
not commutation of physical crossing operators. Edge scalars are
already in $B_N$ and are not counted in $\widehat\zeta_N$.

Choose $R<R'<1$, $R'\ne1/2$, and initially take
$|\sigma_v|=R'$. For one coordinate use the exact pole kernel
\[
 K_\sigma(\tau)=\frac{\ii N e^{\ii N\tau}}
              {e^{\ii N\tau}-\eta e^{2\pi\ii\sigma}}.
\]
It has residue one at precisely the shifted grid points. Integrate
it times $\sin\tau f(\tau)e^{\ii nN\tau^2/(4\pi)}$ around a
rectangle translated by $2\pi\sigma/N$ enclosing those $N$ poles.
On its vertical sides
\[
 K_\sigma(\pm\pi+\ii v+2\pi\sigma/N)
                  =\frac{\ii N e^{-Nv}}{1+e^{-Nv}}.
\]
On horizontal sides expand the denominator geometrically. After
reflecting the upper side by $\tau\mapsto-\tau$, the odd sine
factor pairs the positive and negative frequencies. Their parameter
multiplier is
\[
 \frac{e^{2\pi\ii\ell\sigma}-e^{-2\pi\ii\ell\sigma}}
      {e^{2\pi\ii\sigma}-e^{-2\pi\ii\sigma}}
       =U_{\ell-1}(\cos2\pi\sigma).
\]
The grid sign and the contour orientations give
\begin{equation}\label{ms:charges}
 c_{\ell,N}(\sigma)=\eta^{\ell-1}
                 U_{\ell-1}(\cos2\pi\sigma),\qquad c_{1,N}=1.
\end{equation}
The zero-frequency real integral vanishes by oddness. Its endpoint
currents, and the connectors from the translated endpoints, remain
in the contour identity. This calculation explains why replacing
the repeated-node functional by values would lose the factor
$c_{\ell,N}(0)=\eta^{\ell-1}\ell$.

Choose $t_0>0$ small and put
$L_v=\lfloor t_0|n_v|\rfloor$, $y_v=A_0t_0/|n_v|$.
The finite real functional retained from this exact contour is
\begin{equation}\label{ms:finite-charges}
 \mathcal P_{n,\sigma,N}^{L}f
 =-\frac{N^2 C_{N,n}(\sigma)}{\pi\ii}
 \sum_{\ell=1}^{L}c_{\ell,N}(\sigma)
 \int_{-\pi}^{\pi}\sin\tau f(\tau)
             e^{\ii nN\tau^2/(4\pi)-\ii\ell N\tau}\,d\tau.
\end{equation}
Here is a uniform bound for what was omitted. On the lower
horizontal side the packet and chirp have real rate at most
\[
 F_0-ct^2+\frac{|n||t|y}{2\pi}-\ell y+Cy^2.
\]
The maximum of the two $t$ terms is $n^2y^2/(16\pi^2c)$.
For $\ell>L$ the resulting upper rate is
\[
 F_0+t_0^2\left(\frac{A_0^2}{16\pi^2c}-A_0
                         +\frac{CA_0^2}{n^2}\right)
                         -(\ell-L-1)y.
\]
First choose $A_0$ so its first positive term is at most $A_0/4$,
then a lower threshold for $|n|$ so the last one is also at most
$A_0/4$. The loss is at least $A_0t_0^2/2$, independently of $n$.
The bound $|c_{\ell,N}(\sigma)|\le\ell e^{2\pi R'(\ell-1)}$
makes the remaining geometric sum convergent for all sufficiently
large $N$. Constants depending on fixed $n$ do not change this rate.

At the vertical endpoints the angular loss is fixed. Their chirp
cost is $O(N|n|y)=O(Nt_0)$, and the retained-charge connector cost
is $O(NLy)=O(Nt_0^2)$. Decrease $t_0$ to preserve the loss. This
also controls the zero-mode current. All spectator variables can
remain real or on their physical translated grids: those grid
heights are $O(1/N)$ and their fixed chirps have bounded modulus.
Use the exact telescoping identity
\[
 \prod_v\mathcal D_v-\prod_v\mathcal P_v
 =\sum_j(\prod_{v<j}\mathcal P_v)
                   (\mathcal D_j-\mathcal P_j)
                   (\prod_{v>j}\mathcal D_v).
\]
The other $\mathcal P$'s are finite real integrals of polynomial
operator cost, and the other $\mathcal D$'s are finite physical
grid sums. Thus every telescoping term has rate $F_0-\gamma$ for
one fixed $\gamma>0$. No spectator is moved to a slope-dependent
complex height during another coordinate's tail estimate.

The whole trace-minus-finite-approximation is holomorphic in each
$\sigma_v$. Apply the maximum principle successively to extend
the estimate from the parameter circles to $|\sigma_v|\le R$.
This includes all normal limits at $0,\pm1/2$. For any fixed
upper-half-plane compact use the same translated contour directly;
its sine denominator is bounded away from zero and its endpoint
translation is still $O(1/N)$. Constant vertices can independently
be brought to zero by the preceding disk argument. Fixed
derivatives follow on smaller compacts.

\subsection{A common critical branch for every retained charge}
After truncating the real integral to a small angular box, with
fixed loss, its phase is
\[
 T_{\boldsymbol\ell}(\boldsymbol\tau)
 =\Psi(\boldsymbol\tau)
       +\frac{\ii}{4\pi}\sum_v n_v\tau_v^2
       -\ii\sum_v\ell_v\tau_v.
\]
Set $\varepsilon_v=1/n_v$, $t_v=\ell_v/n_v$. Its critical equation
and a removable stationary expression are
\begin{align}
 \tau_v&=2\pi t_v+2\pi\ii\varepsilon_v\Psi_v(\boldsymbol\tau),
       \label{ms:joint-critical}\\
 W(\mathbf t,\boldsymbol\varepsilon)
 &=\Psi(\boldsymbol\tau)
                -\ii\pi\sum_v\varepsilon_v\Psi_v(\boldsymbol\tau)^2.
\end{align}
The implicit function theorem gives one holomorphic branch on a
fixed polydisk about $(0,0)$. The critical action is
$-\ii\pi\sum t_v^2/\varepsilon_v+W$; the first term is purely
imaginary for the real signed data. Separate evenness makes $W$
even in each $t_v$. At $\boldsymbol\varepsilon=0$,
$W(\mathbf t,0)=\Psi(2\pi\mathbf t)$, so
$\partial_{t_v^2}\Re W(0,0)=-2\pi^2 a_v$.
Shrink the polydisk until all these derivatives are at most
$-\pi^2a_v$. Integration in the squared coordinates gives
\begin{equation}\label{ms:joint-loss}
 \Re T_{\boldsymbol\ell}(\boldsymbol\tau_{\boldsymbol\ell})
 \le\Re T_{\mathbf1}(\boldsymbol\tau_{\mathbf1})
       -\sum_v\frac{\pi^2a_v}{n_v^2}(\ell_v^2-1).
\end{equation}
Choose $t_0$ inside this polydisk and then increase the single
lower threshold for all $|n_v|$. This covers every retained charge
without a restriction on their slope ratios.

Write $\Psi_v=\tau_v Q_v(\tau_1^2,\ldots,\tau_H^2)$, with
$\Re Q_v<0$ in that neighborhood. Then
\begin{equation}\label{ms:critical-angle}
 \tau_{\boldsymbol\ell,v}=
 \frac{2\pi t_v}{1-2\pi\ii\varepsilon_v Q_v},\qquad
 \Im\tau_{\boldsymbol\ell,v}
 =\frac{4\pi^2\ell_v\Re Q_v}
              {n_v^2|1-2\pi\ii\varepsilon_vQ_v|^2}<0.
\end{equation}
Every angle and every sine in the amplitude is nonzero. Translate
the small real box to this critical point. The real part of its
Hessian on real tangent vectors is that of $\Psi''$, since the
chirp Hessian is purely imaginary, and stays negative definite.
The connectors preserve the face loss: their heights are at most
$C|t_v|/|n_v|$, their chirp cost on a face of size $\rho$ is at
most $CN\rho|t_v|$, and their Fourier factors decrease under the
downward displacement. Choose $t_0$ small compared with the fixed
face loss before increasing the slope threshold.

Use the holomorphic packet expansion on this contour, not a real
pointwise remainder. Gaussian integration and
\eqref{ms:dimension} give for the entire retained contraction
\begin{equation}\label{ms:filled-packet}
 \left(\prod_v\mathcal D_v\right)B_N
 =\left(\prod_v C_{N,n_v}(\sigma_v)\right)
   \Gamma_{e,N}(s)N^{3/2}D_{\mathbf n}
   e^{N T_{\mathbf1}(\boldsymbol\tau_{\mathbf1})}
                         (1+O(N^{-1})),
\end{equation}
where
\[
 D_{\mathbf n}=\left(-\frac1{\pi\ii}\right)^H(2\pi)^{H/2}
 \frac{a(\boldsymbol\tau_{\mathbf1})\prod_v\sin\tau_{\mathbf1,v}}
 {\sqrt{\det[-T_{\mathbf1}''(\boldsymbol\tau_{\mathbf1})]}}\ne0.
\]
The branch is fixed by the translated real Gaussian. There are
only finitely many charges for fixed powers. Equation
\eqref{ms:joint-loss} suppresses every one except $\mathbf1$,
whose coefficient is exactly one by \eqref{ms:charges}. The
reflected signs have already been paired, leaving no equal-action
terms with an unexamined relative sign. The wire Hessian and this
angular Schur complement are invertible, hence so is the full
joint Hessian.

The first action tends to $F_0$ as $\max_v|1/n_v|\to0$.
Increase the common threshold to put its real part above $F_0$
minus half each fixed source, contour-tail and off-window loss.
Those additive errors are then negligible. The charge-dependent
losses were instead compared to this same filled action in
\eqref{ms:joint-loss}; a smallest local loss is never compared
with a total complete-to-filled drop.

Finally restrict $\sigma_v$ to the actual types. The finite scalar
ledger \cite[Proposition 8.1]{Meridian}, which retains every typed
turn, monodromy, crossing carrier and framing correction, says
\[
 \frac{\widehat\zeta_N(s)}{\widehat\zeta_N(0)}
 \Gamma_{e,N}(s)\prod_v q^{-n_v\sigma_v(s)^2/2}
                       =q^{a(D)s^2/2},\qquad
 a(D)=-w(D)-\sum_v n_v.
\]
There is still exactly one main-cut factor $H_N$ in
\eqref{ms:source-to-functional}. This proves
\eqref{ms:filled-leading}. Removing these exact scalars leaves
only the product of \eqref{ms:exact-insertion} on the original
physical edges in the joint integral. The zero-meridian amplitude
and all its offsets are common, so no first-order scalar is hidden
in that assertion. Every choice of an exponential margin or a
slope threshold preceded the choice of a meridian compact; only
bounded factors and starting levels depend on that compact.
